\documentclass[10pt,a4paper]{amsart}
\usepackage[margin=19mm]{geometry}
\usepackage{amsmath,amssymb,mathtools}
\usepackage[scr=rsfs,cal=euler]{mathalfa}
\usepackage{xfrac}
\usepackage[unicode,hidelinks,bookmarks=false]{hyperref}
\newcommand{\ie}{i.e.\ }
\newcommand{\cf}{cf.\ }
\newcommand{\resp}{respectively\ }
\newcommand{\Subsection}{\subsection}
\providecommand{\nobreakdash}{\nobreak\hbox{-}}
\setlength{\emergencystretch}{2em}
\multlinegap0pt
\usepackage{amssymb}
\newtheorem{theorem}{Theorem}
\title{A new unknotting operation for classical and welded knots}
\author{Danish Ali, Zhiqing Yang, Mohd Ibrahim Sheikh, Sidra Batool}
\date{Source: arXiv:2208.08761v3 (CC BY 4.0)}
\begin{document}
\maketitle

\noindent\emph{Source and licence.} A new unknotting operation for classical and welded knots. Version arXiv:2208.08761v3 (CC BY 4.0), distributed by arXiv under the Creative Commons Attribution 4.0 International licence, http://creativecommons.org/licenses/by/4.0/, as recorded in the arXiv licence metadata for this exact version and shown on the abstract page https://arxiv.org/abs/2208.08761v3. Full licence notice: \texttt{licenses/arxiv-2208.08761-CC-BY-4.0.txt}.\par\smallskip

\noindent\textit{Editorial scope.} Counted unit: the article's theorem theorem (source0.tex lines 166--168). The article's complete original proof of it is delivered with it (source0.tex lines 170--174, one \texttt{proof} environment of 5 lines). No mathematics is written, repaired or re-derived: the statement and the proof are the article's own lines, in the article's own notation, and no retained line is substituted or re-flowed.  No further source-local context is needed: every cross-reference of every retained line resolves inside the two delivered spans.  Nothing was omitted from any retained span, so the package carries no omission record.  No retained line contains a citation, so the wrapper carries no bibliography and no bibliography entry.  No retained line contains a figure environment, an image-inclusion command or drawing code, so no image is delivered and the package declares no asset dependency.  Editorial formatting only, in addition: an AMS \texttt{amsart} preamble that restates the article's own declarations and macros used by the retained lines, namely the environments \texttt{theorem} on separate counters, together with the standard packages those lines need.  The retained environments print in this document as Theorem 1 in order of appearance, whereas the article prints the same items with the numbering of its own full document.  The complete native source of the article as distributed in the arXiv e-print for this exact version is bundled unchanged as \texttt{sources/129612/source0.tex}.
\begin{theorem}
 Every link diagram can be unlinked by a finite sequence of $D$-moves and Reidemeister moves.
\end{theorem}

\begin{proof}
Let $L = l_1 \cup l_2 \cup \cdots \cup l_m$ be an oriented link diagram, where each component $l_i$ is given a chosen orientation and a base point, and the components are given a fixed order $l_1, l_2, \cdots, l_m$. A crossing of $L$ is either a self-crossing, involving strands of the same component, or a mixed crossing, involving strands from two distinct components. The unlinking process begins by converting the diagram into a completely descending form. Starting with $l_1$, we traverse it once from its base point along its orientation. Whenever we meet a crossing for the first time along $l_i$, if it is a self-crossing and $l_i$ passes under at this first encounter, we perform a crossing change by a $D$-move so that $l_i$ passes over. If it is a mixed crossing with $l_j$, where $j>i$, and $l_i$ is under on the first encounter, we change it so that $l_i$ passes over. If $j<i$, no change is made because that crossing was already encountered while processing $l_j$. This procedure is then applied successively to $l_2, l_3, \dots, l_m$, ensuring that each component always passes over any later component at their first encounter.

By construction, the resulting diagram is completely descending with respect to the chosen base points, orientations, and order of components: for each component, every first encounter with a crossing occurs as an overpass, and for every mixed crossing between $l_i$ and $l_j$ with $i<j$, $l_i$ is over $l_j$. Such a diagram represents the trivial $m$-component unlink. Indeed, one can realize the diagram in $\mathbb{R}^3$ so that $l_1$ lies entirely in a small height interval above the others, $l_2$ lies just below $l_1$, and so on, respecting the over/under information. This spatial separation produces a split union of the components. Within each height slab, the projection of a component is a descending knot diagram, which is known to represent the unknot. Each component can then be isotoped to a round circle in its slab, yielding the standard unlink.
\end{proof} 

\end{document}
